% =====================================================================
% TaskScheduler — Diploma Project
% University Politehnica of Bucharest, Faculty of Automatic Control and
% Computers, Computer Science and Engineering Department.
% =====================================================================
\documentclass[12pt,a4paper]{report}

% ---------- packages -------------------------------------------------
\usepackage[utf8]{inputenc}
\usepackage[T1]{fontenc}
\usepackage[english]{babel}
\usepackage{graphicx}
\usepackage{geometry}
\geometry{a4paper,left=3cm,right=2.5cm,top=2.5cm,bottom=2.5cm}
\usepackage{amsmath,amssymb}
\usepackage{xcolor}
\usepackage{hyperref}
\hypersetup{colorlinks=true,linkcolor=black,citecolor=blue,urlcolor=blue}
\usepackage{listings}
\lstset{
    basicstyle=\ttfamily\footnotesize,
    keywordstyle=\bfseries\color{blue!60!black},
    commentstyle=\itshape\color{green!40!black},
    stringstyle=\color{red!60!black},
    breaklines=true,
    showstringspaces=false,
    frame=single,
    framesep=4pt,
    captionpos=b,
    tabsize=2,
}
\usepackage{booktabs}
\usepackage{array}
\usepackage{tikz}
\usetikzlibrary{positioning,arrows.meta,shapes.geometric,fit,backgrounds}
\usepackage{caption}
\usepackage{subcaption}
\usepackage{enumitem}
\usepackage{setspace}
\onehalfspacing

% ---------- per-thesis metadata --------------------------------------
\newcommand{\Name}{<Your Full Name>}
\newcommand{\Advisor}{<Advisor Name, e.g.\ Prof.\ dr.\ ing.\ XXXX>}
\newcommand{\Year}{2026}
\newcommand{\ProjectTitle}{TaskScheduler: A Heterogeneity- and Thermal-Aware
Workflow Scheduler for Kubernetes}
\newcommand{\ProjectSubtitle}{Two-tier scheduling with predictive ECT,
adaptive UCB exploration, and live bandwidth/thermal telemetry}

% ---------- title page template (provided by the course) -------------
\newcommand{\HeaderLineSpace}{-0.25cm}
\newcommand{\UniTextRO}{UNIVERSITATEA POLITEHNICA DIN BUCUREȘTI \\[\HeaderLineSpace]
FACULTATEA DE AUTOMATICĂ ȘI CALCULATOARE \\[\HeaderLineSpace]
DEPARTAMENTUL DE CALCULATOARE\\}
\newcommand{\DiplomaRO}{PROIECT DE DISERTAȚIE}
\newcommand{\AdvisorRO}{Coordonator științific:}
\newcommand{\BucRO}{BUCUREȘTI}

\newcommand{\UniTextEN}{UNIVERSITY POLITEHNICA OF BUCHAREST \\[\HeaderLineSpace]
FACULTY OF AUTOMATIC CONTROL AND COMPUTERS \\[\HeaderLineSpace]
COMPUTER SCIENCE AND ENGINEERING DEPARTMENT\\}
\newcommand{\DiplomaEN}{DIPLOMA PROJECT}
\newcommand{\AdvisorEN}{Thesis advisor:}
\newcommand{\BucEN}{BUCHAREST}

\newcommand{\frontPage}[6]{
\begin{titlepage}
\begin{center}
{\Large #1}
\vspace{50pt}
\begin{tabular}{p{6cm}p{4cm}}
\includegraphics[scale=0.8]{pics/upb-logo.jpg} &
    \includegraphics[scale=0.5,trim={14cm 11cm 2cm 5cm},clip=true]{pics/cs-logo.pdf}
\end{tabular}

\vspace{105pt}
{\Huge #2}\\
\vspace{40pt}
{\Large #3}\\ \vspace{0pt}
{\Large #4}\\
\vspace{40pt}
{\LARGE \Name}\\
\end{center}
\vspace{60pt}
\begin{tabular*}{\textwidth}{@{\extracolsep{\fill}}p{6cm}r}
&{\large\textbf{#5}}\vspace{10pt}\\
&{\large \Advisor}
\end{tabular*}
\vspace{20pt}
\begin{center}
{\large\textbf{#6}}\\
\vspace{0pt}
{\normalsize \Year}
\end{center}
\end{titlepage}
}

% =====================================================================
\begin{document}

\frontPage{\UniTextEN}{\DiplomaEN}{\ProjectTitle}{\ProjectSubtitle}{\AdvisorEN}{\BucEN}

% ---------- abstract -------------------------------------------------
\chapter*{Abstract}
\addcontentsline{toc}{chapter}{Abstract}

Modern data-processing pipelines run on increasingly heterogeneous
infrastructure: clusters mix CPU-optimised, memory-optimised and
I/O-optimised nodes, run on a variety of CPU architectures (e.g.\
\texttt{amd64} and \texttt{arm64}), and operate under thermal limits
that the default Kubernetes scheduler is unaware of. This project
proposes \emph{TaskScheduler}, a workflow-aware scheduler for
Kubernetes that places tasks using a predictive estimate of effective
completion time (ECT) refined by live measurements of inter-node
bandwidth and per-node CPU temperature, and that explores the
scheduling space using a UCB-style adaptive policy.

TaskScheduler is implemented as a standard Kubernetes control plane:
a CRD-driven controller (\texttt{ts-controller}) materialises workflow
DAGs into Pods, a custom scheduler (\texttt{ts-scheduler}) binds those
Pods to nodes, and three DaemonSets supply the producer-local data
plane (\texttt{ts-fileserver}) and the telemetry signals
(\texttt{ts-bw-probe}, \texttt{ts-thermal}) consumed by the policy.
The system is validated both in a deterministic simulator
(25/25 unit tests) and on a two-host \texttt{k3s} cluster spanning
\texttt{amd64} and \texttt{arm64} hosts, where it can be compared
against the default Kubernetes scheduler under identical workloads.

\vspace{1em}
\noindent\textbf{Keywords:} Kubernetes, workflow scheduling,
heterogeneous clusters, thermal-aware scheduling, bandwidth-aware
scheduling, UCB, ECT, CRDs, k3s.

% ---------- table of contents ----------------------------------------
\tableofcontents

% =====================================================================
\chapter{Introduction}
\label{ch:intro}

\section{Subject and context}
Workflow scheduling on Kubernetes is dominated by the cluster's default
scheduler, which optimises for resource fit and basic affinity
constraints but treats nodes as essentially interchangeable. In
practice, modern clusters are far from uniform: some nodes have many
fast CPU cores but little memory, others the opposite; some sit behind
a passively cooled enclosure and throttle quickly under sustained
load; some are physically far enough from the workflow's data
producer that the network becomes the bottleneck. The default
scheduler has no visibility into any of these signals.

This project addresses that gap for the specific case of
\emph{data-processing pipelines} expressed as directed acyclic graphs
(DAGs) of tasks. The system, called \emph{TaskScheduler}, treats each
workflow as a first-class object (a Kubernetes Custom Resource), tracks
its tasks through their full lifecycle, and binds them to nodes using a
policy that combines structural information (the DAG, node profiles)
with live telemetry (inter-node bandwidth, CPU temperature).

\section{Motivation}
The motivation is twofold.

\textbf{Practical.} Containerised workflows are now the dominant way of
expressing batch and stream pipelines in research and industry
(scientific simulation, ETL, ML training preprocessing, CI), and
Kubernetes is the dominant orchestrator. Yet the default scheduler's
choices are often visibly suboptimal on heterogeneous clusters --- a
memory-heavy task may be placed on the fastest CPU node simply because
it was first to be considered, leaving the memory-optimised node idle
or, worse, occupied by an unrelated CPU task. Even small placement
improvements compound across DAGs with tens or hundreds of tasks.

\textbf{Scientific.} ECT-style estimation, multi-armed-bandit
exploration (UCB), and thermal awareness have each been studied
individually in the scheduling literature, but rarely combined and
almost never integrated cleanly into Kubernetes' control loop. The
present project examines whether the three signals reinforce each
other, and whether the resulting policy beats both the default
Kubernetes scheduler and simpler heuristics (random, round-robin,
best-fit) on realistic mixed workloads.

\section{Objectives}
The project pursues the following concrete objectives:
\begin{enumerate}[label=O\arabic*., leftmargin=*]
    \item Design a workflow-aware scheduling model that captures
    node heterogeneity (CPU/memory/IO profile, cooling class,
    architecture) and DAG structure (upward rank, critical path).
    \item Define an effective-completion-time estimator
    $\widehat{\mathrm{ECT}}(\tau, n \mid A)$ that takes the current
    assignment context $A$ (parent placements, data sizes) into
    account.
    \item Combine ECT with an adaptive UCB exploration policy so the
    scheduler keeps refining its priors during operation rather than
    relying solely on offline profiles.
    \item Integrate live signals --- inter-node bandwidth and per-node
    CPU temperature --- into the scoring function $J(A)$.
    \item Implement the scheduler as a clean Kubernetes deployment
    (CRDs, controller, custom scheduler, DaemonSets) rather than a
    monolithic process, so that it can be deployed on real clusters.
    \item Validate the implementation end-to-end on a two-host
    \texttt{k3s} cluster spanning two CPU architectures, and compare
    it against the default Kubernetes scheduler under identical
    workloads.
\end{enumerate}

\section{Expected impact}
If the proposed policy is shown to consistently reduce workflow
makespan or improve fairness on heterogeneous clusters relative to
the default scheduler --- without requiring intrusive changes to user
workloads --- it offers a directly deployable improvement to a very
widely used piece of cluster software. Even a negative or partial
result is useful: it would clarify which of the three signals (ECT,
UCB, thermal) carry weight in practice and which can be dropped to
simplify the system.

% =====================================================================
\chapter{State of the art (optional)}
\label{ch:sota}

\emph{This chapter is intentionally short; the depth will be expanded
during the final write-up. It briefly situates the work relative to
the most directly relevant prior research.}

\paragraph{Heterogeneous scheduling.} HEFT~\cite{topcuoglu2002heft}
introduced the use of \emph{upward rank} on the DAG combined with an
earliest-finish-time placement rule on heterogeneous processors, and
remains the canonical baseline for workflow schedulers. Many later
variants extend HEFT with stochastic execution times, energy budgets,
or cloud-style pricing.

\paragraph{Bandit-style schedulers.} Multi-armed-bandit and UCB-style
exploration~\cite{auer2002ucb} have been used to learn execution-time
priors online; their appeal is that they give a principled answer to
the exploration/exploitation trade-off without requiring a model of
the cluster.

\paragraph{Thermal-aware scheduling.} Hot nodes throttle, so placing
work on a cool node can finish faster even when the cool node is
nominally slower. Recent work on edge and mobile clusters has shown
measurable gains from incorporating temperature into placement.

\paragraph{Kubernetes-native schedulers.} The Kubernetes scheduler
framework allows plugins to inject scoring logic into the default
binder; alternatively, a custom scheduler can run alongside the
default one, claiming Pods whose \texttt{spec.schedulerName} matches
its own. This project takes the latter approach because it allows
the full ECT + UCB + thermal policy to be expressed cleanly outside
the default scheduler's plugin model.

\paragraph{Workflow controllers.} Argo Workflows and similar tools
demonstrate the CRD-driven controller pattern for DAG execution.
TaskScheduler uses the same pattern, but pairs it with a custom
\emph{scheduler}; the controller is intentionally thin.

% =====================================================================
\chapter{Project model}
\label{ch:model}

This chapter describes the architecture of TaskScheduler, the
scheduling algorithm, and the live signals it consumes.

\section{High-level architecture}
TaskScheduler runs entirely inside Kubernetes as four cooperating
components and three DaemonSets (Figure~\ref{fig:arch}):

\begin{enumerate}[leftmargin=*]
    \item \textbf{Workflow CRD} --- a user submits a workflow as
    a Kubernetes object containing the DAG, the per-task templates
    and the workflow-level metadata (priority, deadline).
    \item \textbf{\texttt{ts-controller} (Deployment).} Watches
    Workflow CRDs, expands the DAG into Pods, sets
    \texttt{schedulerName=ts-scheduler}, scrapes
    \texttt{\_\_TS\_OUTPUT\_\_=} lines from finished pods' logs and
    injects them as environment variables on downstream pods.
    \item \textbf{\texttt{ts-scheduler} (Deployment).} Watches Pending
    Pods whose \texttt{schedulerName} is \texttt{ts-scheduler}, runs
    the policy described in \S\ref{sec:algo}, and binds each Pod to a
    node.
    \item \textbf{Data plane DaemonSet (\texttt{ts-fileserver}).}
    Exposes per-node outputs at \texttt{http://<node>:8080} so that
    downstream tasks running on a different node can fetch parent
    outputs.
    \item \textbf{Telemetry DaemonSets.} \texttt{ts-bw-probe} measures
    inter-node bandwidth every 10\,min and writes
    \texttt{ts.io/bw-to-<peer>=<MB/s>} as Node annotations.
    \texttt{ts-thermal} reads
    \texttt{/sys/class/thermal/thermal\_zone0/temp} every 30\,s and
    writes \texttt{ts.io/cpu-temp-c}.
\end{enumerate}

\begin{figure}[ht]
\centering
\begin{tikzpicture}[
    node distance=8mm,
    box/.style={draw, rounded corners, align=center,
        minimum width=3.2cm, minimum height=1cm, font=\small},
    ds/.style={box, fill=blue!8},
    cp/.style={box, fill=orange!12},
    arr/.style={-{Stealth}, thick}
  ]
    \node[cp] (ctrl) {\texttt{ts-controller}\\(Deployment)};
    \node[cp, right=20mm of ctrl] (sched) {\texttt{ts-scheduler}\\(Deployment)};
    \node[box, above=10mm of $(ctrl)!0.5!(sched)$] (crd) {Workflow CRD\\(user input)};

    \node[ds, below=12mm of ctrl, xshift=-15mm] (fs)   {\texttt{ts-fileserver}\\DaemonSet};
    \node[ds, below=12mm of ctrl, xshift=20mm] (bw)    {\texttt{ts-bw-probe}\\DaemonSet};
    \node[ds, below=12mm of sched, xshift=20mm] (th)   {\texttt{ts-thermal}\\DaemonSet};

    \draw[arr] (crd) -- (ctrl);
    \draw[arr] (ctrl) -- node[above,font=\scriptsize]{Pods} (sched);
    \draw[arr] (sched.south) -- ++(0,-4mm) -| node[pos=0.25,above,font=\scriptsize]{bind} (bw.north);
    \draw[arr] (fs.north) -- node[right,font=\scriptsize]{outputs} (ctrl.south);
    \draw[arr] (bw.east) -- ++(8mm,0) |- node[pos=0.25,right,font=\scriptsize]{bw} (sched.east);
    \draw[arr] (th.north) -- node[right,font=\scriptsize]{temp} (sched.south);
\end{tikzpicture}
\caption{High-level architecture: a thin controller, a custom
scheduler, and three DaemonSets supplying data plane and telemetry.}
\label{fig:arch}
\end{figure}

\section{Two-tier scheduling: from workflows to tasks}
The scheduler operates at two levels of granularity:
\begin{enumerate}
    \item \textbf{Workflow level.} Each workflow accumulates virtual
    runtime $v_w$ (CFS-style); the next workflow to admit a task is
    the one with the smallest $v_w$, ensuring long-term fairness
    across concurrently submitted workflows.
    \item \textbf{Task level.} Within the chosen workflow, the next
    Pending task is the one with the largest \emph{upward rank}
    --- the longest weighted path from that task to any DAG sink ---
    so that critical-path tasks are placed first.
\end{enumerate}
This decoupling is the same idea HEFT uses, generalised to the
multi-workflow case.

\section{Scoring function and ECT}
\label{sec:algo}
Given a candidate (task $\tau$, node $n$) under the partial assignment
$A$ already in place, the scheduler estimates the effective completion
time:
\begin{equation}
\widehat{\mathrm{ECT}}(\tau, n \mid A) =
\max\bigl(t_{\text{ready}}(n),\;
          \max_{p \in \text{parents}(\tau)} t_{\text{done}}(p)
          + \mathrm{transfer}(p, n \mid A)\bigr)
+ \mathrm{exec}(\tau, n),
\label{eq:ect}
\end{equation}
where $\mathrm{transfer}(p, n \mid A)$ uses the live
\texttt{ts.io/bw-to-<peer>} annotation when $p$'s placement is
already decided, and $\mathrm{exec}(\tau, n)$ uses the per-profile
priors maintained by the UCB layer.

The placement objective is then
\begin{equation}
J(A) = \widehat{\mathrm{ECT}}(\tau, n \mid A)
     + \lambda_{\text{therm}} \cdot \mathrm{penalty}_{\text{therm}}(n)
     + \lambda_{\text{fail}}  \cdot \mathrm{penalty}_{\text{fail}}(n)
     - \beta_{\text{ucb}} \cdot \mathrm{bonus}_{\text{ucb}}(\tau, n),
\label{eq:J}
\end{equation}
which the scheduler \emph{minimises} over the set of feasible
candidate nodes. The thermal penalty rises sharply once the node's
recent average CPU temperature crosses an envelope; the UCB bonus
falls as a node accumulates samples for a given task profile.

\section{Live signals}
Two signals are collected continuously by DaemonSets:
\begin{description}[leftmargin=*]
    \item[Inter-node bandwidth.] \texttt{ts-bw-probe} on each node
    fetches the \texttt{\_probe.bin} file (100\,MB) from every peer's
    \texttt{ts-fileserver}, divides size by elapsed time, and writes
    \texttt{ts.io/bw-to-<peer>=<MB/s>} to its own Node object.
    Refreshed every 10\,min. Consumed by
    \texttt{services/bandwidth.py}.
    \item[CPU temperature.] \texttt{ts-thermal} reads
    \texttt{/sys/class/thermal/thermal\_zone0/temp} on Linux and falls
    back to a load-average proxy on hosts without sysfs thermals
    (e.g.\ Multipass VMs on macOS). Refreshed every 30\,s.
    Consumed by \texttt{services/thermal.py} and by the
    \texttt{AdaptivePolicy}.
\end{description}

\section{Producer-local data plane}
Each node owns a host-path directory \texttt{/var/lib/ts-data} that
the controller mounts into every worker Pod at \texttt{/data/outputs}.
Same-node parents are read through that mount with no network
involvement. For cross-node parents, the controller injects
\texttt{TS\_PARENT\_<NAME>\_FILESERVER\_URL} so the task script can
\texttt{curl} the parent's output through the producer's
\texttt{ts-fileserver}. This is intentionally simple: no shared
filesystem, no object store --- the data plane and the bandwidth
signal share exactly the same path, so they stay consistent.

\section{Adaptive policy}
The policy itself is selected via the \texttt{TS\_POLICY} environment
variable on the scheduler Deployment. The default value
\texttt{adaptive} enables the full ECT + UCB + thermal scoring of
Eq.~\ref{eq:J}. Setting \texttt{random}, \texttt{round-robin} or
\texttt{best-fit} switches in simpler baselines, which is what makes
controlled A/B comparisons against the default Kubernetes scheduler
feasible.

% =====================================================================
\chapter{Research methodology}
\label{ch:method}

\section{Implementation as the foundation}
A working implementation is treated as a prerequisite for any
quantitative claim; the simulator and the cluster are not afterthoughts
but the core methodological apparatus. Two execution surfaces are
maintained in parallel:

\begin{enumerate}[leftmargin=*]
    \item A \textbf{discrete-event simulator}
    (\texttt{run\_simulation.py}, \texttt{engine.py}) that runs the
    same policy code path against a synthetic cluster and synthetic
    workflows. It is deterministic, fast (entire batches finish in
    seconds), and is the regression harness: 25 unit tests in
    \texttt{test\_simulation.py} pin the algorithm's invariants.
    \item A \textbf{real Kubernetes cluster}, deployed in two
    configurations:
    \begin{enumerate}[label=(\alph*)]
        \item single-host \texttt{kind} with three worker nodes, used
        for smoke tests and CI;
        \item two-host \texttt{k3s} (described in \S\ref{sec:topo}),
        used for the actual measurements.
    \end{enumerate}
\end{enumerate}

\section{Experimental topology}
\label{sec:topo}
The two-host cluster spans two physically distinct machines on the
same gigabit LAN, providing a real network hop (essential for the
bandwidth signal to mean anything) and two CPU architectures
(essential to confirm the system actually works on heterogeneous
hardware).

\begin{table}[ht]
\centering
\small
\begin{tabular}{lll>{\raggedright\arraybackslash}p{6.2cm}}
\toprule
\textbf{Role} & \textbf{Host} & \textbf{Arch} & \textbf{Logical k8s nodes} \\
\midrule
Master & Mac mini M2, 8\,GB (macOS + Multipass) & arm64 &
\texttt{ts-master} (control plane + scheduler + registry, MEM\_OPT, PASSIVE) \\
Worker & Laptop, Ryzen 7 PRO 8840HS, 64\,GB (Windows + Multipass) & amd64 &
\texttt{ts-node-cpu-1}, \texttt{ts-node-cpu-2} (CPU\_OPT), \texttt{ts-node-io-1} (IO\_OPT), \texttt{ts-node-mem-1} (MEM\_OPT) \\
\bottomrule
\end{tabular}
\caption{Two-host \texttt{k3s} topology. Five logical nodes across two
physical machines and two CPU architectures. The arm64 master runs
inside a Linux VM on macOS; the four amd64 workers run inside Linux
VMs on Windows.}
\label{tab:topo}
\end{table}

The \texttt{ts-scheduler} runs on the master; the Mac mini is purely a
worker, which both reflects a realistic deployment and confirms that
the scheduler does not need to be co-located with the workload.

\section{Workload}
A representative workload mix is used in all experiments. It contains:
\begin{itemize}[leftmargin=*]
    \item \emph{CPU-bound} tasks (\texttt{tasks/task\_cpu}) that
    iterate a tight numeric loop;
    \item \emph{Memory-bound} tasks (\texttt{tasks/task\_mem}) that
    allocate and touch large buffers;
    \item \emph{I/O-bound} tasks (\texttt{tasks/task\_io}) that fetch
    or write large files through the data plane.
\end{itemize}
DAG shapes range from linear pipelines (\texttt{linear-pipeline-1}) to
fan-out/fan-in topologies that exercise cross-node parent-to-child
data transfers.

\section{Comparison protocol}
Every experiment is run twice against the same workload and the same
cluster: once with \texttt{schedulerName=ts-scheduler}
(\texttt{TS\_POLICY=adaptive}), and once with the workflow's Pods
left to the default Kubernetes scheduler. Each pair is repeated
$n=10$ times to obtain confidence intervals. Random seeds and probe
schedules are controlled so the two halves of a pair see the same
cluster state.

\section{Reproducibility}
The full setup is automated:
\texttt{make k3s-server} and \texttt{make k3s-agent} build the cluster
on the two hosts, \texttt{make images-multiarch} cross-builds all
seven container images for \texttt{amd64+arm64} and pushes them to a
LAN registry, and \texttt{make verify-cluster-run} executes the
end-to-end sanity check including a workflow run. All scripts are
idempotent and tear-down is one command (\texttt{make k3s-uninstall}).
The full source --- including this document --- is versioned in Git.

% =====================================================================
\chapter{Contributions}
\label{ch:contrib}

\section{Theoretical contributions}
\begin{enumerate}[leftmargin=*]
    \item A \emph{two-tier} scheduling formulation that combines
    CFS-style workflow-level fairness (vruntime) with HEFT-style
    task-level criticality (upward rank), so that multiple workflows
    can be scheduled concurrently without starvation while each
    workflow still respects its own critical path.
    \item An \emph{ECT estimator}
    $\widehat{\mathrm{ECT}}(\tau, n \mid A)$ (Eq.~\ref{eq:ect}) that
    is computed against the \emph{current} partial assignment $A$,
    using live bandwidth annotations to refine the transfer term and
    UCB-maintained priors for the execution term.
    \item A unified \emph{scoring function} $J(A)$ (Eq.~\ref{eq:J})
    that combines ECT with thermal and failure penalties and a UCB
    exploration bonus.
\end{enumerate}

\section{Practical contributions}
\begin{enumerate}[leftmargin=*]
    \item A complete, production-shaped Kubernetes deployment:
    a Workflow CRD, a controller, a custom scheduler, and three
    DaemonSets, packaged as plain manifests under \texttt{manifests/}.
    The CRDs and the controller/scheduler split mean the system can
    coexist with any other Kubernetes workload on the same cluster.
    \item A \emph{producer-local data plane} that sidesteps the need
    for a shared filesystem: parents write to a host-path,
    \texttt{ts-fileserver} re-exposes that path over HTTP, and the
    controller injects per-parent URLs into child Pods.
    \item Two telemetry signals integrated as Node annotations
    rather than as a separate metrics system, so any Kubernetes
    client can read them with a single \texttt{kubectl get nodes}.
    \item Multi-architecture container images
    (\texttt{linux/amd64} + \texttt{linux/arm64}) built via
    \texttt{docker buildx} and a LAN-local insecure registry,
    enabling the same images to run on a Ryzen laptop and on an
    Apple-silicon Mac mini.
    \item A reproducible experimental harness: \texttt{cluster.env},
    three idempotent setup scripts (\texttt{k3s-server.sh},
    \texttt{k3s-agent.sh}, \texttt{k3s-uninstall.sh}), and a
    \texttt{verify-cluster.sh} health-check that confirms every
    invariant (nodes Ready, labels applied, DaemonSets covering every
    node, bandwidth and thermal annotations populated, an end-to-end
    workflow finishing).
    \item A simulator that runs the \emph{same} policy code on a
    synthetic cluster, giving cheap regression coverage (25/25 unit
    tests) and an easy way to A/B different policies before deploying
    them on the real cluster.
\end{enumerate}

% =====================================================================
\chapter{Performance analysis metrics}
\label{ch:metrics}

The evaluation uses the following metrics, all computed per workflow
and aggregated across runs:

\begin{description}[leftmargin=*]
    \item[Makespan ($M_w$).] Wall-clock time from the first Pod of
    workflow $w$ entering the \texttt{Pending} phase to the last Pod
    of $w$ reaching \texttt{Succeeded}. Primary headline metric.
    \item[Per-task slowdown ($s_\tau$).]
    $s_\tau = T_\tau / T_\tau^{\text{best}}$, where $T_\tau^{\text{best}}$
    is the execution time the same task achieves alone on its
    profile-matched node. Captures how much each task suffered from
    placement choices.
    \item[Critical-path stretch.] Ratio of observed makespan to the
    DAG's critical-path length using profile-matched execution times.
    Decouples policy quality from raw cluster capacity.
    \item[Fairness across workflows.] Jain's fairness index applied
    to per-workflow normalised makespan in multi-workflow runs.
    \item[Thermal exposure.]
    $\mathrm{TE}(n) = \int (\max(0, T_{\mathrm{cpu}}(n,t) - T_{\text{env}}))\,dt$
    integrated over the run. Captures how aggressively the policy
    pushes any single node.
    \item[Bandwidth utilisation.] Mean and 95th percentile of the
    actual data-transfer volume on each inter-node edge, compared
    against the bandwidth measured by \texttt{ts-bw-probe}.
    \item[Cross-node ratio.] Fraction of parent$\to$child edges in
    the DAG for which the child ran on a different node than the
    parent. Lower is better when transfers are large; the policy
    should trade this off against ECT.
    \item[Scheduling latency.] Time from Pod \texttt{Pending} to Pod
    \texttt{Scheduled}. The custom scheduler must not be noticeably
    slower than the default scheduler.
    \item[Failure / retry count.] Number of Pod failures and how
    quickly the scheduler avoids the failing node afterwards.
\end{description}

Each metric is reported as mean $\pm$ 95\,\% CI over the ten repeats
of each (policy, workload) pair, with TaskScheduler and the default
Kubernetes scheduler side by side.

% =====================================================================
\chapter{Results}
\label{ch:results}

\emph{This chapter is intentionally left empty for now; it will be
filled once the planned experiments on the two-host cluster are
complete. The expected structure mirrors the metrics defined in
Chapter~\ref{ch:metrics}: one section per metric, presenting
TaskScheduler against the default Kubernetes scheduler under
identical workloads, with mean and 95\,\% CI and a short discussion of
the observed differences.}

\section{Makespan}
% TODO: fill after experiments

\section{Per-task slowdown and critical-path stretch}
% TODO: fill after experiments

\section{Fairness across concurrent workflows}
% TODO: fill after experiments

\section{Thermal exposure and bandwidth utilisation}
% TODO: fill after experiments

\section{Scheduling latency and overhead}
% TODO: fill after experiments

\section{Discussion}
% TODO: fill after experiments

% =====================================================================
\chapter{Conclusion and further work}
\label{ch:conclusion}

\section{Summary}
This project designed and implemented TaskScheduler, a workflow-aware,
heterogeneity- and thermal-aware scheduler for Kubernetes, and built
the experimental harness needed to evaluate it. The system is
deployable as a regular Kubernetes workload (CRDs, a controller, a
custom scheduler, and three DaemonSets) and has been validated
end-to-end on a two-host \texttt{k3s} cluster spanning \texttt{amd64}
and \texttt{arm64} hardware.

\section{Further research}
Several extensions are natural and were deliberately kept out of scope
of this thesis:
\begin{itemize}[leftmargin=*]
    \item \textbf{Energy- and cost-aware scoring} --- the same $J(A)$
    framework can absorb a per-node energy or cost term once metering
    is available.
    \item \textbf{Preemption} --- the current scheduler does not
    preempt; adding HEFT-style task migration could improve makespan
    under sudden thermal events.
    \item \textbf{Federated clusters} --- exposing the Workflow CRD
    through KubeFed or a similar layer would let a single workflow
    span clusters with different policies.
    \item \textbf{Learned ECT estimators} --- replacing the
    profile-based execution priors with a small online regressor
    that takes task features as input.
\end{itemize}

\section{Closing remarks}
TaskScheduler's main contribution is showing that the three signals
(ECT, UCB, thermal) can be combined cleanly inside a real Kubernetes
control plane, not just in a simulator. The empirical evaluation in
Chapter~\ref{ch:results} will determine how much each of those
signals actually buys in practice.

% =====================================================================
\begin{thebibliography}{9}
\bibitem{topcuoglu2002heft}
H.~Topcuoglu, S.~Hariri, and M.-Y. Wu,
\emph{Performance-effective and low-complexity task scheduling for
heterogeneous computing}, IEEE Trans. on Parallel and Distributed
Systems, 13(3):260--274, 2002.

\bibitem{auer2002ucb}
P.~Auer, N.~Cesa-Bianchi, and P.~Fischer,
\emph{Finite-time analysis of the multiarmed bandit problem},
Machine Learning, 47(2--3):235--256, 2002.

\bibitem{kubernetes}
The Kubernetes Authors, \emph{Kubernetes documentation},
\url{https://kubernetes.io/docs/}, accessed 2026.

\bibitem{k3s}
Rancher Labs, \emph{k3s: Lightweight Kubernetes},
\url{https://k3s.io/}, accessed 2026.
\end{thebibliography}

\end{document}
